\chapter{燃料--包壳间隙传热模型补充关系式}
\label{app:gap_heat_transfer_supplementary_models}


% 为兼容原有子公式交叉引用：分组后保留旧标签，并统一指向所在关系组的公式编号。
\makeatletter
\providecommand{\appendixeqlabelalias}[1]{%
  \begingroup
  \edef\@currentlabel{\p@equation\theequation}%
  \label{#1}%
  \endgroup}
\makeatother

本附录汇总正文\ref{sec:gap_heat_transfer_model}中未展开的几何修正、温度跳跃、热适应系数、混合气体热导率、高压气体物性和接触传导关系。相关模型主要由轻水堆$\mathrm{UO}_2$/锆合金燃料棒试验与BISON、MATPRO实现发展而来；用于快堆MOX/不锈钢体系时，应分别核查几何、气体状态及界面材料参数，不应直接继承全部默认值。

\section{等效间隙长度与几何修正}
\label{app:gap_geometry_models}

Ross--Stoute关系中的$d_g$为与传热面积定义一致的等效间隙长度。对于平板、圆柱和球形间隙，BISON采用\cite{halesBISONTheoryManual2016,toptanGapConductanceModeling2019}
\begin{equation}
	\label{eq:appendix_gap_effective_lengths}
	\begin{aligned}
		d_g&=R_2-R_1,
		&&\text{平板近似}\\
		d_g&=R_{\mathrm{ref}}
		\ln\!\left(\frac{R_2}{R_1}\right),
		&&\text{圆柱间隙}\\
		d_g&=R_{\mathrm{ref}}^2
		\left(\frac{1}{R_1}-\frac{1}{R_2}\right),
		&&\text{球形间隙}
	\end{aligned}
\end{equation}
\appendixeqlabelalias{eq:appendix_gap_plate_length}
\appendixeqlabelalias{eq:appendix_gap_cylinder_length}
\appendixeqlabelalias{eq:appendix_gap_sphere_length}
式中，$R_1$和$R_2$分别为内、外表面半径，$R_{\mathrm{ref}}$为热流密度所对应参考表面的半径。对于$d=R_2-R_1\ll R_1$的薄间隙，三种表达均趋近于几何间隙$d$。有限元实现中，$d_g$必须与界面热流所采用的参考面积保持一致，否则会重复引入曲率修正。

\section{温度跳跃与热适应系数}
\label{app:gap_temperature_jump_models}

\subsection{Lanning--Hann模型}

Lanning--Hann模型基于Kennard理论，将两侧表面的总温度跳跃距离写为\cite{lanningReviewMethodsApplicable1975,halesBISONTheoryManual2016}
\begin{equation}
	g_f+g_c
	=C_L
	\left(\frac{2-\alpha_{\mathrm{mix}}}
	{\alpha_{\mathrm{mix}}}\right)
	\frac{k_g\sqrt{T_g}}{p_g}
	\left(\sum_{i=1}^{N_g}\frac{x_i}{M_i}\right)^{-1/2}
	\label{eq:appendix_gap_lanning_jump_distance}
\end{equation}
其中，$\alpha_{\mathrm{mix}}$为混合气体热适应系数。BISON的传统实现取$C_L=5756$，该常数对应以下旧制单位：
\begin{itemize}
	\item $g_f+g_c$采用$\mathrm{cm}$，$k_g$采用$\mathrm{cal\,cm^{-1}\,s^{-1}\,K^{-1}}$；
	\item $p_g$采用$\mathrm{dyn\,cm^{-2}}$，$M_i$采用$\mathrm{g\,mol^{-1}}$。
\end{itemize}
若程序内部统一使用SI单位，应先对式\eqref{eq:appendix_gap_lanning_jump_distance}进行整体单位换算，不得只替换单个输入量。

MATPRO传统模型对氦和氙给出\cite{hagrmanSCDAPRELAP5MOD1995}
\begin{equation}
	\label{eq:appendix_gap_legacy_accommodation}
	\begin{aligned}
		\alpha_{\mathrm{He}}
		&=0.425-2.3\times10^{-4}T_g\\
		\alpha_{\mathrm{Xe}}
		&=0.749-2.5\times10^{-4}T_g\\
		\alpha_{\mathrm{mix}}
		&=\alpha_{\mathrm{He}}
		+\frac{\left(\alpha_{\mathrm{Xe}}-\alpha_{\mathrm{He}}\right)
		\left(M_{\mathrm{mix}}-M_{\mathrm{He}}\right)}
		{M_{\mathrm{Xe}}-M_{\mathrm{He}}}
	\end{aligned}
\end{equation}
\appendixeqlabelalias{eq:appendix_gap_he_accommodation}
\appendixeqlabelalias{eq:appendix_gap_xe_accommodation}
\appendixeqlabelalias{eq:appendix_gap_mixture_accommodation}
该插值主要面向He--Xe混合物。数值实现应将$\alpha$限制在$0<\alpha\leq1$的物理范围内；含Ar、Kr或其他组分时，应采用与实际气体--表面组合相匹配的多组分模型。

\subsection{Toptan模型}

Toptan等在曲线坐标下重新整理了Kennard温度跳跃关系，并显式引入两侧表面热适应系数及几何修正\cite{toptanGapConductanceModeling2019,toptanModelingGapConductance2020}：
\begin{equation}
	g_T
	=\frac{C_K k_g}{p_g\Lambda_{12}}
	\sqrt{\frac{T_g}{R_gM_g}}
	\label{eq:appendix_gap_toptan_jump_distance}
\end{equation}
其中，$R_g$为气体常数，$M_g$为气体摩尔质量，$C_K$为与分子自由度有关的Kennard系数。BISON采用
\begin{equation}
	C_K=
	\begin{cases}
	0.2173,&\text{单原子气体}\\
	0.1149,&\text{双原子气体}\\
	0.1242,&\text{多原子气体}
	\end{cases}
	\label{eq:appendix_gap_kennard_coefficients}
\end{equation}
两表面的几何--适应系数为
\begin{equation}
	\label{eq:appendix_gap_lambda_geometry}
	\begin{aligned}
		\Lambda_{12}^{\mathrm{plate}}
		&=\frac{\alpha_1\alpha_2}
		{\alpha_1+\alpha_2-\alpha_1\alpha_2}\\
		\Lambda_{12}^{\mathrm{cyl}}
		&=\frac{\alpha_1\alpha_2}
		{\alpha_2+\alpha_1(1-\alpha_2)(R_1/R_2)}\\
		\Lambda_{12}^{\mathrm{sph}}
		&=\frac{\alpha_1\alpha_2}
		{\alpha_2+\alpha_1(1-\alpha_2)(R_1/R_2)^2}
	\end{aligned}
\end{equation}
\appendixeqlabelalias{eq:appendix_gap_lambda_plate}
\appendixeqlabelalias{eq:appendix_gap_lambda_cylinder}
\appendixeqlabelalias{eq:appendix_gap_lambda_sphere}

Toptan模型采用经实验数据标定的热适应系数关系
\begin{equation}
	\label{eq:appendix_gap_toptan_accommodation}
	\begin{aligned}
		\alpha(T_g)
		&=1-
		\left\{
		1-\alpha_{\infty}
		\tanh\!\left[
		\frac{\sqrt{M_gT_g}}{\alpha_{\infty}}\theta_1
		\right]
		\right\}
		\exp\!\left(-\frac{\theta_2}{T_g}\right)\\
		\alpha_{\infty}
		&=\frac{2.4\mu}{(1+\mu)^2},
		\qquad
		\mu=\frac{M_g}{M_s}
	\end{aligned}
\end{equation}
\appendixeqlabelalias{eq:appendix_gap_toptan_alpha}
\appendixeqlabelalias{eq:appendix_gap_toptan_alpha_infinity}
标定中固体表观摩尔质量取$M_s=10^5\ \mathrm{g\,mol^{-1}}$。系数见表\ref{tab:appendix_gap_toptan_accommodation_parameters}。

\begin{table}[htbp]
	\centering
	\small
	\caption{Toptan热适应系数模型参数\cite{toptanModelingGapConductance2020}}
	\label{tab:appendix_gap_toptan_accommodation_parameters}
	
	\begin{tabular*}{0.80\textwidth}{
			@{\extracolsep{\fill}}
			cccccc
			@{}
		}
		\toprule
		参数 & He & Ne & Ar & Kr & Xe \\
		\midrule
		$\theta_1$ & $-1.055$ & $-0.611$ & $0.032$ & $0.344$ & $0.784$ \\
		$\theta_2$ & $146.7$ & $360.6$ & $591.3$ & $724.7$ & $916.7$ \\
		\bottomrule
	\end{tabular*}
\end{table}

\section{混合气体热导率}

\subsection{MATPRO温度相关模型}
\label{app:gap_default_gas_conductivity}

正文式\eqref{eq:gap_gas_mixture_conductivity}中的二元权重系数可写为\cite{brokawApproximateFormulasViscosity1958,hagrmanSCDAPRELAP5MOD1995}
\begin{equation}
	\label{eq:appendix_gap_brokaw_weights}
	\begin{aligned}
		\Psi_{ij}
		&=\phi_{ij}
		\left[
		1+2.41
		\frac{(M_i-M_j)(M_i-0.142M_j)}{(M_i+M_j)^2}
		\right]\\
		\phi_{ij}
		&=\frac{
		\left[
		1+\left(k_i/k_j\right)^{1/2}
		\left(M_i/M_j\right)^{1/4}
		\right]^2}
		{2^{3/2}\left(1+M_i/M_j\right)^{1/2}}
	\end{aligned}
\end{equation}
\appendixeqlabelalias{eq:appendix_gap_psi_ij}
\appendixeqlabelalias{eq:appendix_gap_phi_ij}
纯气体热导率采用
\begin{equation}
	k_i(T_g)=A_iT_g^{B_i}
	\label{eq:appendix_gap_pure_gas_conductivity}
\end{equation}
其中，$k_i$的单位为$\mathrm{W\,m^{-1}\,K^{-1}}$。MATPRO及高温补充数据对应的参数见表\ref{tab:appendix_gap_gas_conductivity_coefficients}\cite{hagrmanSCDAPRELAP5MOD1995,springerThermalConductivityNeon1973}。

\begin{table}[htbp]
	\centering
	\small
	\caption{纯气体热导率关系$k_i=A_iT^{B_i}$的经验系数}
	\label{tab:appendix_gap_gas_conductivity_coefficients}
	\begin{tabular*}{0.80\textwidth}{
			@{\extracolsep{\fill}}
			lccc
			@{}
		}
		\toprule
		气体 & $A_i$ & $B_i$ & 温度范围（K） \\
		\midrule
		He  & $2.639\times10^{-3}$ & 0.7085 & 73--793 \\
		Ne  & $9.000\times10^{-4}$ & 0.6910 & 73--793 \\
		Ne  & $9.683\times10^{-4}$ & 0.6850 & 1000--1500 \\
		Ar  & $2.986\times10^{-4}$ & 0.7224 & 73--793 \\
		Ar  & $4.905\times10^{-4}$ & 0.6510 & 1000--2500 \\
		Kr  & $8.247\times10^{-5}$ & 0.8363 & 173--793 \\
		Xe  & $4.351\times10^{-5}$ & 0.8616 & 173--793 \\
		Xe  & $1.140\times10^{-4}$ & 0.7100 & 1000--1500 \\
		$\mathrm{H}_2$ & $1.097\times10^{-3}$ & 0.8785 & -- \\
		$\mathrm{N}_2$ & $5.314\times10^{-4}$ & 0.6898 & -- \\
		$\mathrm{O}_2$ & $1.853\times10^{-4}$ & 0.8729 & -- \\
		CO & $1.403\times10^{-4}$ & 0.9090 & -- \\
		$\mathrm{CO}_2$ & $9.460\times10^{-6}$ & 1.3120 & -- \\
		\bottomrule
	\end{tabular*}
\end{table}

若事故工况需要考虑水蒸气进入燃料棒自由体积，可采用MATPRO关系
\begin{equation}
	k_{\mathrm{H_2O}}
	=17.6\times10^{-3}
	+T_C\left[
	5.87\times10^{-5}
	+T_C\left(1.04\times10^{-7}
	-4.51\times10^{-11}T_C\right)
	\right]
	\label{eq:appendix_gap_steam_conductivity}
\end{equation}
其中，$T_C=T_g-273.15$，$T_C$的单位为$^\circ\mathrm C$。对于正常密封的快堆燃料棒，水蒸气通常不属于初始或裂变生成气体，式\eqref{eq:appendix_gap_steam_conductivity}默认未启用。

\subsection{温度--压力相关高级模型}
\label{app:gap_advanced_gas_conductivity}

MATPRO幂函数主要依据低压数据标定。在极低压力的Knudsen区，气体热导率随压力近似增加；在约$1\ \mathrm{MPa}$以下的常用工程区间，压力影响通常较弱；压力进一步升高后，密度效应可能不可忽略\cite{toptanGapConductanceModeling2019,tournierPropertiesNobleGases2008}。Tournier--El-Genk模型将纯惰性气体热导率写为
\begin{equation}
	\label{eq:appendix_gap_advanced_conductivity_model}
	\begin{aligned}
		k_i(T_g,p_g)
		&=k_i^0(T_g)+k_{i,c}^{*}\Psi_k(\rho_{r,i})\\
		k_i^0(T_g)
		&=\frac{15R_g}{4M_i}\mu_i^0(T_g)\\
		\mu_i^0(T_g)
		&=A_{\mu,i}\left(T_g-T_{\mu,i}\right)^{n_i}\\
		k_{i,c}^{*}
		&=\frac{0.201\times10^{-4}
		T_{c,i}^{0.277}M_i^{-0.465}}
		{\left(0.291V_i^{*}\right)^{0.415}},
		\qquad
		V_i^{*}=\frac{R_gT_{c,i}}{p_{c,i}}\\
		\Psi_k(\rho_{r,i})
		&=0.645\rho_{r,i}
		+0.331\rho_{r,i}^2
		+0.0368\rho_{r,i}^3
		-0.0128\rho_{r,i}^4,
		\qquad
		\rho_{r,i}=\frac{\rho_i}{\rho_{c,i}}
	\end{aligned}
\end{equation}
\appendixeqlabelalias{eq:appendix_gap_advanced_k}
\appendixeqlabelalias{eq:appendix_gap_dilute_k}
\appendixeqlabelalias{eq:appendix_gap_dilute_viscosity}
\appendixeqlabelalias{eq:appendix_gap_pseudocritical_k}
\appendixeqlabelalias{eq:appendix_gap_excess_conductivity}
所需临界参数见表\ref{tab:appendix_gap_critical_gas_properties}。

\begin{table}[htbp]
	\centering
	\small
	\caption{高级气体热导率模型的临界参数\cite{tournierPropertiesNobleGases2008}}
	\label{tab:appendix_gap_critical_gas_properties}
	\begin{tabular*}{0.90\textwidth}{
			@{\extracolsep{\fill}}
			lccccc
			@{}
		}
		\toprule
		参数 & He & Ne & Ar & Kr & Xe \\
		\midrule
		$M$（$\mathrm{g\,mol^{-1}}$） & 4.003 & 20.180 & 39.948 & 83.798 & 131.293 \\
		$p_c$（MPa） & 0.2275 & 2.678 & 4.863 & 5.51 & 5.84 \\
		$T_c$（K） & 5.2 & 44.5 & 150.69 & 209.4 & 289.7 \\
		$\rho_c$（$\mathrm{kg\,m^{-3}}$） & 69.64 & 481.9 & 535.6 & 908.4 & 1110.0 \\
		$A_\mu\times10^7$ & 3.063 & 8.4528 & 6.989 & 6.963 & 7.568 \\
		$T_\mu$（K） & $-21.33$ & 16.47 & 65.70 & 71.07 & 112.31 \\
		$n$ & 0.724 & 0.643 & 0.640 & 0.667 & 0.655 \\
		\bottomrule
	\end{tabular*}
\end{table}

高级模型通过三阶维里状态方程计算摩尔密度$\hat\rho$：
\begin{equation}
	\frac{p_g}{R_gT_g}
	=\hat\rho+B_2(T_g)\hat\rho^2+B_3(T_g)\hat\rho^3
	\label{eq:appendix_gap_virial_equation}
\end{equation}
其中，$B_2$和$B_3$分别为第二、第三维里系数。BISON在求解式\eqref{eq:appendix_gap_virial_equation}时选取满足模型约束的最大正实根。该处理仅在启用温度--压力相关气体导热模型时需要；对于常压附近的传统MATPRO模型，不应额外重复施加密度修正。

\section{固--固接触传导补充模型}
\label{app:gap_contact_supplementary_models}

\subsection{收缩热阻与真实接触面积}

设$A_0$为表观接触面积、$\dot Q$为通过界面的总热流率，则界面接触热阻和单位面积热阻分别为
\begin{equation}
	\label{eq:appendix_gap_contact_resistance}
	\begin{aligned}
		R_c&=\frac{\Delta T_c}{\dot Q}
		=\frac{1}{h_sA_0}\\
		R_c''&=A_0R_c=\frac{1}{h_s}
	\end{aligned}
\end{equation}
\appendixeqlabelalias{eq:appendix_gap_total_contact_resistance}
\appendixeqlabelalias{eq:appendix_gap_area_contact_resistance}
对于半无限固体之间半径为$a$的单个圆形接触斑点，收缩热阻可近似写为
\begin{equation}
	R_{c,a}=\frac{1}{2ak_m},
	\qquad
	k_m=\frac{2k_fk_c}{k_f+k_c}
	\label{eq:appendix_gap_single_spot_resistance}
\end{equation}
真实接触面积$A_r$与表观面积$A_0$的关系通常表示为
\begin{equation}
	\frac{A_r}{A_0}
	=C_A\left(\frac{P_c}{H}\right)^n
	\label{eq:appendix_gap_real_contact_area}
\end{equation}
其中，$C_A$为表面形貌系数，$n$反映微凸体由弹性向塑性变形的偏离程度。Ross--Stoute模型以粗糙度和接触压力替代难以直接测量的接触斑点数量与尺度，得到正文式\eqref{eq:gap_solid_contact_model}所示工程关联。

\subsection{迈耶硬度的历史关系}

MATPRO/BISON提供的锆合金历史关系为\cite{hagrmanSCDAPRELAP5MOD1995,halesBISONTheoryManual2016}
\begin{equation}
	\begin{aligned}
	H_{\mathrm{Zr}}(T)
	=\exp\!\bigl(&26.034
	-2.6394\times10^{-2}T\\
	&+4.3502\times10^{-5}T^2
	-2.5621\times10^{-8}T^3\bigr)
	\end{aligned}
	\label{eq:appendix_gap_zircaloy_meyer_hardness}
\end{equation}
其中，$H_{\mathrm{Zr}}$的单位为Pa，$T$的单位为K。该式基于锆合金数据。

该关系不适用于SS316、D9或HT9。快堆不锈钢包壳计算应采用对应材料及加工状态下的温度相关硬度；若缺少数据，应将$H$和$C_s$纳入敏感性或不确定性分析，而不应将BISON默认$680\ \mathrm{MPa}$视为材料常数。

